\begin{table}[h]
\centering
\caption{Synthetic-to-empirical bid-coefficient ratios (Study 2).}
\label{tab:slope-ratios-Study2}
\begin{threeparttable}
\footnotesize
\begin{tabular}{lcccc}
\toprule
 & \multicolumn{2}{c}{Random-cost} & \multicolumn{2}{c}{Bounded} \\
\cmidrule(lr){2-3} \cmidrule(lr){4-5}
Model & w/o beliefs & w/ beliefs & w/o beliefs & w/ beliefs \\
\midrule
AGGREGATED &  0.895 &  0.321 &  1.633 & -0.695 \\
\midrule
deepseek-chat-v3.1 &  0.777 &  0.368 &  1.536 & -0.541 \\
deepseek-r1 &  0.866 &  0.402 &  1.987 & -0.394 \\
gemini-2.5-flash &  1.122 &  0.119 &  2.516 & -0.712 \\
gemini-2.5-flash-lite &  1.021 &  0.083 &  2.768 & -0.675 \\
gpt-5-mini &  1.052 &  0.150 &  2.853 & -0.854 \\
kimi-k2 &  1.097 &  0.789 &  3.253 &  0.576 \\
llama-4-scout &  1.071 &  0.442 &  2.811 & -0.396 \\
mistral-medium-3.1 &  1.180 &  0.503 &  1.981 & -0.642 \\
mistral-small-3.2-24b &  1.344 &  0.138 &  1.383 & -1.257 \\
\bottomrule
\end{tabular}
\begin{tablenotes}\footnotesize
\item \textit{Notes:} Ratio of synthetic to empirical bid coefficients from conditional logit estimates. 
Values below 1 indicate that synthetic responses are less sensitive to cost than empirical responses (). 
A value of 0 indicates no sensitivity to cost; values above 1 indicate over-sensitivity. 
The empirical bid coefficients used as denominators are taken from the unadjusted conditional logit specifications reported in Table~\ref{tab:repl-clogit-Study2}.
\end{tablenotes}
\end{threeparttable}
\end{table}
